% ─────────────────────────────────────────────────────────────────────────────
%  OpenISAC – Extensión ZMQ Software-in-the-Loop
%  Documento técnico: arquitectura, modelo de canal, cambios al código original
%  y manual de uso paso a paso
% ─────────────────────────────────────────────────────────────────────────────
\documentclass[11pt,a4paper]{article}

\usepackage{geometry}
\geometry{top=2.5cm, bottom=2.5cm, left=2.8cm, right=2.8cm}

\usepackage{fontspec}
\setmainfont{DejaVu Serif}
\setmonofont{DejaVu Sans Mono}

\usepackage{xcolor}
\usepackage{listings}
\usepackage{tcolorbox}
\tcbuselibrary{listings,skins,breakable}
\usepackage{hyperref}
\usepackage{booktabs}
\usepackage{tabularx}
\usepackage{longtable}
\usepackage{graphicx}
\usepackage{amsmath}
\usepackage{amssymb}
\usepackage{enumitem}
\usepackage{fancyhdr}
\usepackage{titlesec}
\usepackage{parskip}
\usepackage{array}

% ── Colores ──────────────────────────────────────────────────────────────────
\definecolor{primary}{RGB}{0,82,147}
\definecolor{secondary}{RGB}{230,240,255}
\definecolor{codebg}{RGB}{248,248,248}
\definecolor{codeborder}{RGB}{200,210,220}
\definecolor{codegreen}{RGB}{40,140,40}
\definecolor{codegray}{RGB}{100,100,100}
\definecolor{codeblue}{RGB}{0,0,180}
\definecolor{warnbg}{RGB}{255,248,220}
\definecolor{warnborder}{RGB}{200,160,0}
\definecolor{infobg}{RGB}{230,245,255}
\definecolor{infoborder}{RGB}{0,120,200}
\definecolor{mathbg}{RGB}{240,255,240}
\definecolor{mathborder}{RGB}{30,160,30}

% ── Estilo de código ──────────────────────────────────────────────────────────
\lstdefinestyle{cpp}{
  language=C++,
  basicstyle=\scriptsize\ttfamily,
  keywordstyle=\color{codeblue}\bfseries,
  commentstyle=\color{codegreen}\itshape,
  stringstyle=\color{primary},
  numberstyle=\tiny\color{codegray},
  backgroundcolor=\color{codebg},
  frame=single,
  framesep=4pt,
  rulecolor=\color{codeborder},
  breaklines=true,
  numbers=left,
  numbersep=8pt,
  tabsize=4,
  showstringspaces=false,
  captionpos=b,
  xleftmargin=1em,
}

\lstdefinestyle{bash}{
  language=bash,
  basicstyle=\footnotesize\ttfamily,
  keywordstyle=\color{codeblue}\bfseries,
  commentstyle=\color{codegreen}\itshape,
  backgroundcolor=\color{codebg},
  frame=single,
  framesep=4pt,
  rulecolor=\color{codeborder},
  breaklines=true,
  showstringspaces=false,
}

\lstdefinestyle{plain}{
  basicstyle=\footnotesize\ttfamily,
  backgroundcolor=\color{codebg},
  frame=single,
  framesep=4pt,
  rulecolor=\color{codeborder},
  breaklines=true,
  showstringspaces=false,
}

% ── Cajas informativas ────────────────────────────────────────────────────────
\newtcolorbox{infobox}[1][]{
  colback=infobg, colframe=infoborder,
  fonttitle=\bfseries, title=#1,
  breakable, arc=3pt, boxrule=0.8pt
}
\newtcolorbox{warnbox}[1][]{
  colback=warnbg, colframe=warnborder,
  fonttitle=\bfseries, title=#1,
  breakable, arc=3pt, boxrule=0.8pt
}
\newtcolorbox{mathbox}[1][]{
  colback=mathbg, colframe=mathborder,
  fonttitle=\bfseries, title=#1,
  breakable, arc=3pt, boxrule=0.8pt
}

% ── Cabeceras ────────────────────────────────────────────────────────────────
\pagestyle{fancy}
\fancyhf{}
\fancyhead[L]{\small\color{primary}\textbf{OpenISAC — Extensión ZMQ Software-in-the-Loop}}
\fancyhead[R]{\small\color{primary}\thepage}
\fancyfoot[C]{\small\color{codegray}Documento técnico — fork\_OpenISAC · Mayo 2026}
\renewcommand{\headrulewidth}{0.4pt}

\titleformat{\section}{\large\bfseries\color{primary}}{}{0em}{\thesection\quad}
\titleformat{\subsection}{\normalsize\bfseries\color{primary!80}}{}{0em}{\thesubsection\quad}
\titleformat{\subsubsection}{\normalsize\bfseries}{}{0em}{\thesubsubsection\quad}

% ─────────────────────────────────────────────────────────────────────────────
\begin{document}

% ── Portada ───────────────────────────────────────────────────────────────────
\begin{titlepage}
\centering
\vspace*{1.5cm}
{\color{primary}\rule{\linewidth}{2pt}}\vspace{0.4cm}

{\Huge\bfseries\color{primary} OpenISAC}\\[0.4cm]
{\LARGE Extensión ZMQ Software-in-the-Loop}\\[0.3cm]
{\large Canal ISAC simulado con interfaz TUI interactiva}\\[0.3cm]
{\normalsize Documentación técnica: arquitectura, modelo de canal, cambios al código y manual de uso}

\vspace{0.4cm}
{\color{primary}\rule{\linewidth}{2pt}}

\vspace{1.5cm}

\begin{tcolorbox}[colback=secondary, colframe=primary, arc=5pt, boxrule=1pt]
\centering
\textbf{Resumen ejecutivo}\\[0.5em]
\begin{flushleft}
Este documento describe la extensión \textit{software-in-the-loop} (SIL) desarrollada sobre
el proyecto de código abierto OpenISAC (disponible en GitHub). La extensión añade un modo de
simulación via ZeroMQ (\texttt{--zmq}) que permite ejecutar la cadena completa
TX $\rightarrow$ Canal $\rightarrow$ RX de procesado ISAC sin ningún hardware de radio (USRP),
sustituyendo la interfaz SDR por dos herramientas de simulación:

\begin{itemize}[leftmargin=1em,topsep=0pt]
  \item \textbf{ChannelSimulator\_void}: canal transparente (pasa la señal sin modificar),
        útil como referencia para verificar que no exista actividad cuando no hay objetos.
  \item \textbf{ChannelSimulator}: canal radar monostático completo con retardo de rango,
        desplazamiento Doppler, micro-Doppler, ecuación del radar y ruido AWGN, controlable
        en tiempo real mediante una interfaz TUI (ncurses).
\end{itemize}

El documento incluye la formulación matemática rigurosa del modelo de canal, las referencias
exactas al código que implementa cada efecto, y las instrucciones paso a paso para ejecutar
el sistema tanto en modo hardware (USRP) como en modo simulado.
\end{flushleft}
\end{tcolorbox}

\vspace{1.5cm}

\begin{tabular}{ll}
\toprule
\textbf{Proyecto base}     & OpenISAC (fork) — \url{https://github.com/OpenISAC} \\
\textbf{Lenguaje}          & C++17 \\
\textbf{Build system}      & CMake $\geq$ 3.18 \\
\textbf{Plataforma}        & Linux (Ubuntu 22.04+), kernel $\geq$ 5.15 \\
\textbf{Dependencias SIL}  & ZeroMQ $\geq$ 4.3, cppzmq, ncurses, FFTW3 \\
\textbf{Fecha}             & Mayo 2026 \\
\bottomrule
\end{tabular}

\vfill
{\small\color{codegray}
Documento de referencia para publicación en congreso sobre ISAC.
Generado con XeLaTeX — fork\_OpenISAC.
}
\end{titlepage}

\tableofcontents
\newpage

% ─────────────────────────────────────────────────────────────────────────────
\section{Introducción y motivación}

OpenISAC~\footnote{\url{https://github.com/OpenISAC}} es una implementación de código abierto en C++ de un sistema
ISAC (\textit{Integrated Sensing and Communications}) basado en OFDM
(\textit{Orthogonal Frequency-Division Multiplexing}), diseñado para ejecutarse sobre
hardware SDR (\textit{Software-Defined Radio}) de la familia USRP de Ettus Research /
National Instruments.

El proyecto original asume disponibilidad de hardware físico para \emph{cualquier}
prueba del sistema, lo que impone restricciones importantes en entornos de desarrollo,
validación de algoritmos y docencia:
\begin{itemize}
  \item La cadena de procesado ISAC completa no puede ejecutarse sin un USRP conectado.
  \item No existe mecanismo para controlar los parámetros del canal (distancia, velocidad,
        SNR) de forma reproducible entre experimentos.
  \item La validación de los algoritmos de procesado radar (mapa Rango-Doppler, CFAR)
        queda condicionada a la disponibilidad del entorno de medidas.
\end{itemize}

Esta extensión, denominada \textbf{ZMQ SIL} (\textit{Software-in-the-Loop}), resuelve
estos problemas añadiendo:
\begin{enumerate}
  \item Un modo de transmisión ZMQ (\texttt{--zmq}) que sustituye la interfaz USRP por
        sockets ZeroMQ PUSH/PULL, sin modificar la cadena de procesado.
  \item Un canal simulado transparente (\texttt{ChannelSimulator\_void}) para obtener una
        referencia sin efectos de canal.
  \item Un canal simulado físico (\texttt{ChannelSimulator}) que implementa un modelo
        radar monostático parametrizable en tiempo real.
\end{enumerate}

% ─────────────────────────────────────────────────────────────────────────────
\section{Arquitectura del sistema completo}

\subsection{Modos de operación}

El sistema puede operar en dos modos mutuamente excluyentes:

\begin{center}
\begin{tabular}{lll}
\toprule
\textbf{Modo} & \textbf{Interfaz TX/RX} & \textbf{Casos de uso} \\
\midrule
Hardware (normal) & USRP via UHD & Experimentos OTA, medidas reales \\
ZMQ SIL (\texttt{--zmq}) & ZeroMQ sockets & Desarrollo, validación, docencia \\
\bottomrule
\end{tabular}
\end{center}

\subsection{Pipeline de señal en modo ZMQ SIL}

\begin{center}
\begin{tcolorbox}[colback=codebg, colframe=codeborder, arc=3pt, boxrule=0.8pt,
                  title={\small\textbf{Figura 1 — Pipeline ZMQ SIL completo}}]
\small
% --- Fila 1: bloques TX y canal ---
\begin{center}
\begin{tabular}{ccc}
\textbf{OFDMModulator} \texttt{--zmq} & & \textbf{ChannelSimulator} \\[4pt]
\fbox{\begin{tabular}{c}LDPC encode\\OFDM IFFT\\CP insert\\(sin USRP)\end{tabular}}
& $\xrightarrow{\quad\text{ZMQ PUSH :5558}\quad}$
& \fbox{\begin{tabular}{c}Retardo $\tau$\\Doppler $f_d$\\Micro-Doppler\\Ec. Radar\\AWGN $w[n]$\end{tabular}}
\end{tabular}
\end{center}
\vspace{4pt}
\begin{center}
$\Big\downarrow$ \small{ZMQ PUSH :5559}
\end{center}
\vspace{2pt}
% --- Fila 2: bloque RX sensing y visualización ---
\begin{center}
\begin{tabular}{ccc}
\textbf{Sensing RX} \texttt{(en Modulator)} & & \textbf{Visualización} \\[4pt]
\fbox{\begin{tabular}{c}PairedFrameQueue\\OFDM demod\\Range-Doppler FFT\\CFAR\end{tabular}}
& $\xrightarrow{\quad\text{UDP :8888}\quad}$
& \fbox{\begin{tabular}{c}\texttt{plot\_sensing\_fast.py}\\Mapa R-D\\en tiempo real\end{tabular}}
\end{tabular}
\end{center}
\end{tcolorbox}
\end{center}

En modo hardware, la señal \textbf{se transmite por el aire} entre antenas TX y RX del USRP.
En modo ZMQ SIL, la señal se transfiere por \textbf{memoria local} a través de sockets ZeroMQ,
y el canal físico es emulado por software.

\subsection{Diagramas de flujo por modo de ejecución}

Los tres modos de ejecución posibles se ilustran a continuación de forma independiente.

\subsubsection{Modo 1 — Hardware real (USRP)}

\begin{center}
\begin{tcolorbox}[colback=codebg, colframe=primary, arc=3pt, boxrule=0.8pt,
                  title={\small\textbf{Figura 2 — Modo Hardware: ejecución normal con USRP}}]
\small\centering
\begin{tabular}{ccccc}
\textbf{Terminal 1} & & \textbf{Canal físico} & & \textbf{Terminal 2} \\[4pt]
\fbox{\begin{tabular}{c}OFDMModulator\\(sin --zmq)\\LDPC + OFDM\\UHD TX stream\end{tabular}}
& $\xrightarrow{\text{aire (RF)}}$
& \fbox{\begin{tabular}{c}USRP\\TX/RX\\antenas reales\\propagación\end{tabular}}
& $\xrightarrow{\text{UHD RX}}$
& \fbox{\begin{tabular}{c}SensingChannel\\Range-Doppler\\UDP :8888\end{tabular}} \\[6pt]
\multicolumn{5}{c}{$\downarrow$ UDP} \\[2pt]
\multicolumn{5}{c}{\fbox{\texttt{plot\_sensing\_fast.py} — Terminal 3}} \\
\end{tabular}
\vspace{2pt}
\begin{flushleft}
\textbf{Comandos:}\\
\texttt{Terminal 1:} \texttt{cd build \&\& ./OFDMModulator Modulator.yaml}\\
\texttt{Terminal 2:} \texttt{python3 ../scripts/plot\_sensing\_fast.py}
\end{flushleft}
\end{tcolorbox}
\end{center}

\subsubsection{Modo 2 — Canal void (referencia nula, sin USRP)}

\begin{center}
\begin{tcolorbox}[colback=codebg, colframe=codegreen, arc=3pt, boxrule=0.8pt,
                  title={\small\textbf{Figura 3 — Modo Canal Void: sin efectos de canal}}]
\small\centering
\begin{tabular}{ccc}
\textbf{Terminal 1} & & \textbf{Terminal 2} \\[4pt]
\fbox{\begin{tabular}{c}ChannelSimulator\_void\\BIND :5558 (PULL)\\BIND :5559 (PUSH)\\y[n] = x[n]\end{tabular}}
& $\xleftarrow{\text{ZMQ :5558}}$
& \fbox{\begin{tabular}{c}OFDMModulator --zmq\\LDPC + OFDM\\ZMQ PUSH :5558\end{tabular}} \\[6pt]
\multicolumn{3}{c}{$\downarrow$ ZMQ :5559 (eco sin cambios)} \\[4pt]
\multicolumn{3}{c}{\fbox{\begin{tabular}{c}SensingChannel (en Modulator)\\Range-Doppler $\rightarrow$ UDP :8888\end{tabular}}} \\[4pt]
\multicolumn{3}{c}{\fbox{\texttt{plot\_sensing\_fast.py} — Terminal 3 (sin picos de target)}} \\
\end{tabular}
\vspace{2pt}
\begin{flushleft}
\textbf{Comandos (en orden):}\\
\texttt{Terminal 1:} \texttt{cd build \&\& ./ChannelSimulator\_void Modulator.yaml}\\
\texttt{Terminal 2:} \texttt{cd build \&\& ./OFDMModulator --zmq Modulator.yaml}\\
\texttt{Terminal 3:} \texttt{python3 ../scripts/plot\_sensing\_fast.py}
\end{flushleft}
\end{tcolorbox}
\end{center}

\subsubsection{Modo 3 — Canal simulado con efectos físicos (sin USRP)}

\begin{center}
\begin{tcolorbox}[colback=codebg, colframe=warnborder, arc=3pt, boxrule=0.8pt,
                  title={\small\textbf{Figura 4 — Modo Canal Simulado: efectos físicos controlables}}]
\small\centering
\begin{tabular}{ccc}
\textbf{Terminal 1} & & \textbf{Terminal 2} \\[4pt]
\fbox{\begin{tabular}{c}ChannelSimulator\\TUI ncurses\\Retardo $\tau(R)$\\Doppler $f_d(v)$\\Micro-Doppler $(A_m,f_m)$\\Ec. Radar $A_\text{eco}(R,\sigma)$\\AWGN $w[n]$\end{tabular}}
& $\xleftarrow{\text{ZMQ :5558}}$
& \fbox{\begin{tabular}{c}OFDMModulator --zmq\\LDPC + OFDM\\ZMQ PUSH :5558\end{tabular}} \\[6pt]
\multicolumn{3}{c}{$\downarrow$ ZMQ :5559 (señal con efectos del canal)} \\[4pt]
\multicolumn{3}{c}{\fbox{\begin{tabular}{c}SensingChannel (en Modulator)\\Range-Doppler $\rightarrow$ picos en $(R_i, v_i)$\\UDP :8888\end{tabular}}} \\[4pt]
\multicolumn{3}{c}{\fbox{\texttt{plot\_sensing\_fast.py} — Terminal 3 (picos visibles en el mapa R-D)}} \\
\end{tabular}
\vspace{2pt}
\begin{flushleft}
\textbf{Comandos (en orden):}\\
\texttt{Terminal 1:} \texttt{cd build \&\& ./ChannelSimulator Modulator.yaml}\\
\texttt{Terminal 2:} \texttt{cd build \&\& ./OFDMModulator --zmq Modulator.yaml}\\
\texttt{Terminal 3:} \texttt{python3 ../scripts/plot\_sensing\_fast.py}
\end{flushleft}
\end{tcolorbox}
\end{center}

\subsection{Flujo de datos detallado}

\begin{enumerate}
  \item \textbf{OFDMModulator (hilo \_modulation\_proc)}:
        codifica los datos de usuario con LDPC, modula con OFDM (IFFT + CP) y deposita
        cada frame de $N_\text{sym} \times (N_{FFT} + N_{CP})$ muestras complejas en
        el buffer circular (\texttt{\_circular\_buffer}).

  \item \textbf{OFDMModulator (hilo \_tx\_proc)} — modo ZMQ:
        extrae el frame del buffer circular y lo envía via ZMQ PUSH al puerto 5558.
        Simultáneamente, encola los símbolos en frecuencia en \texttt{PairedFrameQueue}
        (lado TX) para el procesado radar posterior.

  \item \textbf{ChannelSimulator}:
        recibe el frame por ZMQ PULL (puerto 5558), aplica los efectos del canal
        (véase Sección~\ref{sec:canal}) y envía el resultado por ZMQ PUSH (puerto 5559).

  \item \textbf{OFDMModulator (hilo \_zmq\_sensing\_rx\_proc)}:
        recibe el frame procesado por ZMQ PULL (puerto 5559) y lo inyecta en
        \texttt{PairedFrameQueue} (lado RX), haciendo corresponder cada frame RX
        con su frame TX por número de secuencia.

  \item \textbf{OFDMModulator (hilo \_sensing\_loop)}:
        extrae pares (TX, RX) de \texttt{PairedFrameQueue}, aplica demodulación OFDM
        al lado RX y calcula el mapa Rango-Doppler mediante FFT 2D. El resultado
        se envía via UDP al puerto 8888 para visualización con \texttt{plot\_sensing\_fast.py}.
\end{enumerate}

% ─────────────────────────────────────────────────────────────────────────────
\section{Modelo de canal implementado}
\label{sec:canal}

El \texttt{ChannelSimulator} implementa un modelo físico de canal radar monostático
para $N_T \leq 3$ objetivos (\textit{targets}) independientes. El modelo es continuo
en tiempo discreto: se aplica muestra a muestra sobre el flujo I/Q recibido.

La señal de salida del canal es la \textbf{superposición de los ecos de todos los
targets más ruido AWGN}:

\begin{mathbox}[Señal compuesta de salida del canal]
\begin{equation}
  y[n] = \sum_{i=1}^{N_T} A_{\text{eco},i} \cdot x[n - d_i] \cdot p_i[n] + w[n]
  \label{eq:canal_completo}
\end{equation}
donde $x[n]$ es la señal transmitida, $A_{\text{eco},i}$ la amplitud del eco del
target $i$, $d_i$ el retardo en muestras, $p_i[n]$ el fasor de Doppler (con
micro-Doppler opcional) y $w[n]$ el ruido aditivo.
\end{mathbox}

\subsection{Retardo de rango}

El eco del target $i$ a distancia $R_i$ se recibe con un retardo de ida y vuelta:
\begin{equation}
  \tau_i = \frac{2 R_i}{c} \quad \text{[s]}
\end{equation}
donde $c = 2{,}998 \times 10^8$ m/s es la velocidad de la luz. El retardo en muestras
enteras es:
\begin{equation}
  d_i = \left\lfloor \frac{2 R_i}{c} \cdot f_s + 0{,}5 \right\rfloor \quad \text{[muestras]}
  \label{eq:delay}
\end{equation}
con $f_s$ la frecuencia de muestreo del sistema (parámetro \texttt{sample\_rate} en el YAML).

El retardo se implementa mediante un \textbf{buffer circular} de profundidad
$d_\text{max} = \lceil 2 \times 20000 / c \times f_s \rceil + N_\text{frame} + 1$
(suficiente para alojar el eco más lejano de la TUI: $R_\text{max} = 10$ km,
round-trip $= 20$ km):

\begin{lstlisting}[style=cpp,
  caption={Implementación del buffer circular de retardo — ChannelSimulator.cpp, líneas 68--72 y 113--115}]
// Profundidad del buffer circular: round-trip de 20 km + margen
const size_t max_delay_samps =
    static_cast<size_t>(std::ceil(2.0f * 20000.0f / C_LIGHT * fs))
    + buffer_size + 1u;
std::vector<std::complex<float>> ring(max_delay_samps, {0.f, 0.f});

// ...
// Escritura en el buffer circular con el frame entrante
for (ssize_t i = 0; i < n; ++i)
    ring[(ring_write + static_cast<size_t>(i)) % max_delay_samps] = in_buf[i];
\end{lstlisting}

\begin{lstlisting}[style=cpp,
  caption={Cálculo del retardo y lectura desde el buffer — ChannelSimulator.cpp, líneas 129--130 y 156--159}]
// d_i = round(2*R_i / c * f_s)  [ecuacion eq:delay]
const size_t delay_samps =
    static_cast<size_t>(std::round(2.0f * dist / C_LIGHT * fs));

// ...
// Lectura del buffer con el retardo correspondiente al target i
const size_t src =
    (ring_write + static_cast<size_t>(i) + max_delay_samps - delay_samps)
    % max_delay_samps;
\end{lstlisting}

\begin{infobox}[Resolución en rango]
La resolución en rango del procesado ISAC es $\Delta R = c / (2B)$, donde $B$ es el
ancho de banda de la señal OFDM ($B \approx f_s$ para subportadoras densas). Con
$f_s = 50$ MHz: $\Delta R = 3$ m/muestra (muestras de tiempo) vs.
$\Delta R_\text{ISAC} = c/2B \approx 3$ m (bin Range-FFT). El retardo del canal
está cuantificado a resolución de 1 muestra, coherente con la resolución del
procesador ISAC.
\end{infobox}

\subsection{Desplazamiento Doppler}

Un target que se mueve con velocidad radial $v_i$ (positivo = aproximándose)
introduce un desplazamiento Doppler en la señal recibida:
\begin{equation}
  f_{d,i} = \frac{2 v_i f_c}{c} \quad \text{[Hz]}
  \label{eq:doppler_freq}
\end{equation}
donde $f_c$ es la frecuencia portadora (parámetro \texttt{center\_freq} en el YAML).

El desplazamiento Doppler se implementa como una \textbf{multiplicación por un fasor
complejo giratorio de amplitud unitaria}. El fasor se actualiza de forma incremental
entre muestras consecutivas, evitando el cálculo de funciones trigonométricas en el
bucle interno (RT critical path):
\begin{equation}
  \Delta\phi_i = \frac{2\pi f_{d,i}}{f_s}
  \label{eq:doppler_phi}
\end{equation}
\begin{equation}
  p_i[n+1] = p_i[n] \cdot e^{j\Delta\phi_i} = p_i[n] \cdot (\cos\Delta\phi_i + j\sin\Delta\phi_i)
  \label{eq:doppler_phasor}
\end{equation}

La multiplicación del fasor es una operación de \textbf{2 multiplicaciones + 2 sumas reales}
(rotación compleja), frente a las $\approx 15$ ns de \texttt{sin}/\texttt{cos} en hardware
x86. El fasor se renormaliza a módulo unidad cada 256 frames para compensar la deriva
acumulada de punto flotante.

\begin{lstlisting}[style=cpp,
  caption={Inicialización del fasor Doppler — ChannelSimulator.cpp, líneas 146--147}]
// Delta de fase Doppler por muestra: d_phi = 2*pi * f_d / fs  [ec. eq:doppler_phi]
// f_d = 2*vel*fc/c                                           [ec. eq:doppler_freq]
const float   d_phi    = 2.0f * PIf * (2.0f * vel * fc / C_LIGHT) / fs;
const std::complex<float> d_rot{std::cos(d_phi), std::sin(d_phi)};
\end{lstlisting}

\begin{lstlisting}[style=cpp,
  caption={Avance incremental del fasor — ChannelSimulator.cpp, líneas 168--172}]
// Rotacion incremental: p[n+1] = p[n] * d_rot  [ec. eq:doppler_phasor]
// Implementada como multiplicacion de complejos (sin llamar a sin/cos por muestra)
dop_phasor[ti] = {
    dop_phasor[ti].real() * d_rot.real() - dop_phasor[ti].imag() * d_rot.imag(),
    dop_phasor[ti].real() * d_rot.imag() + dop_phasor[ti].imag() * d_rot.real()
};
\end{lstlisting}

\begin{lstlisting}[style=cpp,
  caption={Renormalización del fasor cada 256 frames — ChannelSimulator.cpp, líneas 184--191}]
// Renormalizar phasors cada 256 frames para evitar deriva de punto flotante.
// La multiplicacion acumulada de floats introduce error O(eps * n_samples).
if (++renorm_ctr >= 256) {
    renorm_ctr = 0;
    for (int ti = 0; ti < 3; ++ti) {
        const float mag = std::abs(dop_phasor[ti]);
        if (mag > 0.0f) dop_phasor[ti] /= mag;
    }
}
\end{lstlisting}

\subsection{Micro-Doppler}

El micro-Doppler modela la modulación adicional en fase introducida por la vibración
o movimiento periódico de partes del target (p.ej. hélices, extremidades). Se implementa
como una modulación en fase sinusoidal superpuesta al fasor Doppler principal:

\begin{equation}
  \phi_m(n) = A_m \sin\!\left(\frac{2\pi f_m n}{f_s}\right)
  \label{eq:microdoppler_phase}
\end{equation}

El fasor combinado (Doppler + micro-Doppler) es:
\begin{equation}
  p_i[n] = e^{j\left(\phi_{d,i}(n) + A_m \sin(2\pi f_m n / f_s)\right)}
  \label{eq:phasor_combined}
\end{equation}

donde $A_m$ [\text{rad}] es la amplitud de la desviación de fase y $f_m$ [Hz] la frecuencia
de vibración. Este efecto produce \textit{sidebands} simétricas en el espectro Doppler
separadas $\pm k f_m$ de la frecuencia Doppler principal, característica de objetos
con partes en movimiento.

La implementación avanza el acumulador de micro-Doppler sample a sample mediante las
fórmulas de adición trigonométrica (sin llamadas a \texttt{sin}/\texttt{cos} en el bucle):

\begin{lstlisting}[style=cpp,
  caption={Inicialización del acumulador micro-Doppler — ChannelSimulator.cpp, líneas 149--154}]
// md_phase[ti] acumula la fase del oscilador de micro-Doppler entre frames
// md_step = 2*pi*f_m / fs  (incremento de fase por muestra)
const float md_step   = 2.0f * PIf * md_freq / fs;
float md_sin = std::sin(md_phase[ti]);   // sin(phi_m[n])
float md_cos = std::cos(md_phase[ti]);   // cos(phi_m[n])
const float md_sin_s  = std::sin(md_step);  // sin(Delta_phi_m)
const float md_cos_s  = std::cos(md_step);  // cos(Delta_phi_m)
\end{lstlisting}

\begin{lstlisting}[style=cpp,
  caption={Aplicación del micro-Doppler al fasor — ChannelSimulator.cpp, líneas 161--178}]
std::complex<float> phasor = dop_phasor[ti];
if (md_amp > 0.0f) {
    // phi_micro = A_m * sin(phi_m[n])       [ec. eq:microdoppler_phase]
    // phasor_total = phasor_doppler * exp(j * phi_micro)   [ec. eq:phasor_combined]
    const float md_mod = md_amp * md_sin;
    phasor *= std::complex<float>{std::cos(md_mod), std::sin(md_mod)};
}
out_buf[i] += echo_amp * ring[src] * phasor;  // suma contribucion del target i

// Avance del oscilador de micro-Doppler usando suma de angulos (sin trig por muestra):
// sin(phi + Delta) = sin(phi)*cos(Delta) + cos(phi)*sin(Delta)
// cos(phi + Delta) = cos(phi)*cos(Delta) - sin(phi)*sin(Delta)
if (md_amp > 0.0f) {
    const float ns = md_sin * md_cos_s + md_cos * md_sin_s;
    const float nc = md_cos * md_cos_s - md_sin * md_sin_s;
    md_sin = ns; md_cos = nc;
}
\end{lstlisting}

\subsection{Amplitud del eco — ecuación del radar normalizada}

La potencia de la señal recibida en un radar monostático viene dada por la
\textbf{ecuación del radar}:
\begin{equation}
  P_r = \frac{P_t G_t G_r \lambda^2 \sigma_\text{RCS}}{(4\pi)^3 R^4}
  \label{eq:radar_equation}
\end{equation}
donde $P_t$ es la potencia transmitida, $G_t$ y $G_r$ las ganancias de antena TX y RX,
$\lambda = c/f_c$ la longitud de onda, $\sigma_\text{RCS}$ la sección eficaz radar
(\textit{Radar Cross Section}) y $R$ la distancia al target.

En el dominio de amplitud de señal ($A = \sqrt{P}$), la amplitud del eco recibido escala
como $R^{-2}$. Para la simulación se usa una forma normalizada de la ecuación:

\begin{mathbox}[Amplitud del eco — forma implementada]
\begin{equation}
  A_{\text{eco},i} = K_\text{sim} \cdot G(\theta_i) \cdot \lambda \cdot \frac{\sqrt{\sigma_{\text{RCS},i}}}{R_i^2}
  \cdot 10^{P_{TX}/20} \cdot 10^{G_{RX}/20}
  \label{eq:echo_amp}
\end{equation}

donde $K_\text{sim} = 120$ es la constante de normalización calibrada para obtener
$A_\text{eco} \approx 0.01$ a las condiciones de referencia: $R = 25$ m,
$\sigma_\text{RCS} = 1$ m$^2$, antena omni ($G = 1$), $f_c = 5{,}8$ GHz,
$P_{TX} = G_{RX} = 0$ dB.

\medskip
\textbf{Verificación numérica de $K_\text{sim} = 120$:}

A las condiciones de referencia, $\lambda = c/f_c = 3\!\times\!10^8\,\text{m/s}\;/\;5{,}8\!\times\!10^9\,\text{Hz} \approx 0{,}0517$ m.
Con $G(\theta)=1$ (omni), $\sqrt{\sigma_\text{RCS}}=1$ m, $R^2 = 25^2 = 625$ m² y $P_{TX}=G_{RX}=0$ dB ($10^{0/20}=1$):

\begin{equation}
  A_\text{eco} = K_\text{sim}\cdot \frac{G \cdot \lambda \cdot \sqrt{\sigma_\text{RCS}}}{R^2}
               = 120 \cdot \frac{1 \cdot 0{,}0517 \cdot 1}{625}
               = \frac{6{,}204}{625} \approx 0{,}00993 \approx 0{,}01
\end{equation}

El valor $K_\text{sim} = 120$ (\texttt{KSIM\_NORM} en \texttt{ChannelSimulator.cpp}) se elige de modo que la amplitud del eco sea del orden de $10^{-2}$, lo que sitúa la señal útil entre el piso de ruido ($\sigma_n \approx 10^{-3}$ con $N_\text{dB}=-60$ dBP) y la saturación de la escala unitaria, proporcionando un margen dinámico operativo de $\approx 20$ dB.
\end{mathbox}

\begin{lstlisting}[style=cpp,
  caption={Cálculo de la amplitud del eco — ChannelSimulator.cpp, líneas 105--143}]
// Ganancia TX y RX en amplitud: 10^(dB/20)
const float tx_scale = std::pow(10.0f, tx_power_db / 20.0f);
const float rx_scale = std::pow(10.0f, rx_gain_db  / 20.0f);

// Ganancia de antena maxima (antena directiva Gaussiana)
// G_max = 16*ln2 / theta_3dB^2  (antena de haz Gaussiano)
// Para antena omni: G = 1.
const float G = (ant_type == 0 || ant_bw_deg <= 0.0f) ? 1.0f : [&] {
    const float bw = ant_bw_deg * PIf / 180.0f;     // radianes
    return 16.0f * 0.6931472f / (bw * bw);          // G_max
}();

// ...
// Patron de radiacion Gaussiano para la direccion del target
float beam_factor = 1.0f;
if (ant_type == 1 && ant_bw_deg > 0.0f) {
    const float theta_half = (ant_bw_deg * 0.5f) * PIf / 180.0f;
    const float theta      = std::abs(angle_d) * PIf / 180.0f;
    // G(theta) = exp(-ln2 * (theta / theta_half)^2)   [ec. eq:beam_pattern]
    beam_factor = std::exp(-0.6931472f * (theta / theta_half) * (theta / theta_half));
}

// Amplitud del eco: A_eco = Ksim * G * G(theta) * lambda * sqrt(RCS) / R^2
//                           * tx_scale * rx_scale          [ec. eq:echo_amp]
// lambda = C_LIGHT / fc  (calculado al inicio de signal_loop)
const float echo_amp = KSIM_NORM * G * beam_factor * lambda * std::sqrt(rcs)
                       / (dist * dist) * tx_scale * rx_scale;
\end{lstlisting}

\subsection{Diagrama de radiación de antena directiva}

Para el modo de antena directiva (\texttt{antenna\_type = 1}), el diagrama de radiación
se aproxima con una forma Gaussiana, que es una buena aproximación al lóbulo principal
de antenas de apertura y de bocina:

\begin{equation}
  G(\theta) = G_\text{max} \cdot \exp\!\left(-\ln 2 \cdot \left(\frac{\theta}{\theta_{3\text{dB}}/2}\right)^{\!2}\right)
  \label{eq:beam_pattern}
\end{equation}

donde $\theta$ es el ángulo off-boresight del target, $\theta_{3\text{dB}}$ el ancho de
haz a $-3$ dB (\texttt{antenna\_beamwidth\_deg}) y $G_\text{max}$ la ganancia de pico:
\begin{equation}
  G_\text{max} = \frac{16 \ln 2}{\theta_{3\text{dB}}^2} \quad (\theta_{3\text{dB}} \text{ en radianes})
  \label{eq:gain_max}
\end{equation}

Esta expresión se obtiene imponiendo que $G(\theta_{3\text{dB}}/2) = G_\text{max}/2$
(definición de los 3 dB) en la forma Gaussiana y que la integral del patrón sobre el
espacio sólido total sea la unidad (conservación de energía en la aproximación de campo
lejano).

\begin{infobox}[Para antena omnidireccional]
Cuando \texttt{antenna\_type = 0} (Omni), se usa $G(\theta) = 1$ para todo $\theta$,
lo que equivale a $G_\text{max} = 1$ y $\theta_{3\text{dB}} \to \infty$.
\end{infobox}

\subsection{Ruido AWGN}

Se añade ruido gaussiano blanco aditivo complejo (\textit{Additive White Gaussian Noise})
con densidad espectral de potencia uniforme. El parámetro de control es la \textbf{potencia
del ruido} $N_\text{dB}$ expresada en dB con referencia a potencia unidad:
\begin{equation}
  \sigma_n^2 = 10^{N_\text{dB}/10}, \qquad
  \sigma_n   = 10^{N_\text{dB}/20}
  \label{eq:awgn_sigma}
\end{equation}
\begin{equation}
  w[n] \sim \mathcal{CN}(0, \sigma_n^2), \quad
  w[n] = \sigma_n \cdot \mathcal{N}(0,1) + j \sigma_n \cdot \mathcal{N}(0,1)
  \label{eq:awgn}
\end{equation}

\begin{warnbox}[Nota sobre el modelo de ruido]
El parámetro \texttt{noise\_power\_db} actúa como \textbf{piso de ruido absoluto}:
es independiente de la ganancia de recepción \texttt{rx\_gain\_db}. Esto modela un
sistema donde el ruido entra en la antena (ruido de antena + temperatura ambiente),
mientras que \texttt{rx\_gain\_db} actúa como un amplificador que mejora la SNR del eco
pero no afecta al piso de ruido (figura de ruido nula del amplificador). Para simular
una figura de ruido finita, incrementar \texttt{noise\_power\_db} y \texttt{rx\_gain\_db}
simultáneamente en la misma cantidad.
\end{warnbox}

La generación de ruido usa una \textbf{tabla precalculada} de $2^{17} = 131\,072$ muestras
Gaussianas (512 KB, cabe en caché L3) para evitar la latencia del generador
\texttt{std::normal\_distribution} en el bucle RT ($\approx 15$ ns/muestra vs.
$\approx 2$ ns con tabla):

\begin{lstlisting}[style=cpp,
  caption={Generación de ruido AWGN con tabla precalculada — ChannelSimulator.cpp, líneas 78--86 y 196--206}]
// Tabla de ruido precalculada (una vez al inicio, semilla fija = 42 para reproducibilidad)
std::vector<float> noise_tbl(NOISE_TBL_SZ);   // NOISE_TBL_SZ = 2^17 = 131072
{
    std::mt19937 rng0{42u};
    std::normal_distribution<float> g{0.f, 1.f};
    for (auto& x : noise_tbl) x = g(rng0);
}
size_t nidx = 0;   // indice ciclico

// --- Bucle por muestra (hot path RT) ---
// sigma_n = sqrt(10^(noise_pdb/10))  [ec. eq:awgn_sigma]
const float nsigma = std::sqrt(std::pow(10.0f, noise_pdb / 10.0f));
for (ssize_t i = 0; i < n; ++i) {
    // w[n] = sigma_n * (N_I + j*N_Q)  [ec. eq:awgn]
    // N_I, N_Q ~ N(0,1) independientes, tomados de la tabla ciclica
    out_buf[i] += std::complex<float>{
        nsigma * noise_tbl[nidx        & (NOISE_TBL_SZ - 1u)],   // componente I
        nsigma * noise_tbl[(nidx + 1u) & (NOISE_TBL_SZ - 1u)]    // componente Q
    };
    nidx += 2u;
}
\end{lstlisting}

\subsection{Tabla de parámetros del canal}

{\small
\noindent Todos los campos pertenecen a \texttt{ChannelParams.hpp}.

\begin{tabularx}{\linewidth}{l p{4.6cm} p{2.0cm} X}
\toprule
\textbf{Parámetro} & \textbf{Campo} & \textbf{Defecto} & \textbf{Descripción} \\
\midrule
\multicolumn{4}{l}{\textit{Parámetros de sistema}} \\
Potencia TX & \texttt{tx\_power\_db} & 0 dB & Escala TX: $10^{P_{TX}/20}$ \\
Ganancia RX & \texttt{rx\_gain\_db} & 0 dB & Escala RX: $10^{G_{RX}/20}$ \\
Ruido AWGN  & \texttt{noise\_power\_db} & $-60$ dBP & Piso ruido: $\sigma_n = 10^{N/20}$ \\
Tipo antena & \texttt{antenna\_type} & 0 (Omni) & 0 = omni, 1 = directiva Gauss. \\
Ancho haz   & \texttt{antenna\_beamwidth\_deg} & 30° & $\theta_{3\text{dB}}$ para directiva \\
$N_T$       & \texttt{num\_targets} & 1 & Núm.\ targets activos (0–3) \\
\midrule
\multicolumn{4}{l}{\textit{Parámetros por target (×3)}} \\
$R_i$            & \texttt{distance}          & 25 m   & Distancia al target \\
$v_i$            & \texttt{velocity}          & 10 m/s & Velocidad radial ($+$: acercándose) \\
$\theta_i$       & \texttt{angle\_deg}        & 0°     & Ángulo off-boresight (directiva) \\
$A_m$            & \texttt{micro\_doppler\_amp}  & 0 rad  & Amplitud modulación micro-Doppler \\
$f_m$            & \texttt{micro\_doppler\_freq} & 10 Hz  & Frecuencia vibración micro-Doppler \\
$\sigma_\text{RCS}$ & \texttt{rcs}            & 1 m$^2$ & Sección eficaz radar del target \\
\bottomrule
\end{tabularx}
}

% ─────────────────────────────────────────────────────────────────────────────
\section{Canal transparente: ChannelSimulator\_void}

El \texttt{ChannelSimulator\_void} implementa un canal completamente transparente:
recibe cada frame I/Q por ZMQ PULL (puerto 5558) y lo retransmite \textbf{sin ninguna
modificación} por ZMQ PUSH (puerto 5559):

\begin{equation}
  y[n] = x[n] \quad \forall n
  \label{eq:void_channel}
\end{equation}

Su propósito es proporcionar una \textbf{referencia experimental} que permite:
\begin{itemize}
  \item Verificar que el pipeline ZMQ completo funciona correctamente.
  \item Comprobar que, sin objetos en el canal, el mapa Rango-Doppler no muestra
        picos espurios significativos (solo la correlación cruzada TX/TX cerca del
        origen, que corresponde a la señal directa).
  \item Depurar el procesado radar aislando los efectos del canal.
\end{itemize}

\begin{lstlisting}[style=cpp,
  caption={Bucle principal del canal void — channel\_sim/src/ChannelSimulator\_void.cpp, líneas 30--49}]
// Canal transparente: recv() -> send() sin modificacion
// y[n] = x[n]  [ec. eq:void_channel]
static void signal_loop(RadioInterface& radio, size_t buffer_size) {
    std::vector<std::complex<float>> buf(buffer_size);
    uint64_t frames = 0;

    while (g_running.load(std::memory_order_relaxed)) {
        const ssize_t n = radio.recv(buf.data(), buffer_size);
        if (n <= 0) continue;   // timeout o error — reintentar

        radio.send(buf.data(), static_cast<size_t>(n));  // retransmitir intacto
        ++frames;

        // Indicador de actividad: imprimir cada 1000 frames
        if (frames % 1000u == 0u)
            std::printf("[VOID] Tramas retransmitidas: %llu\n",
                        static_cast<unsigned long long>(frames));
    }
}
\end{lstlisting}

\begin{infobox}[Resultado esperado con canal void]
Al ejecutar \texttt{plot\_sensing\_fast.py} con el canal void activo:
\begin{itemize}
  \item El mapa Rango-Doppler no debe mostrar picos de alta amplitud fuera del
        bin (0,0) (correlación directa).
  \item La presencia del bin (0,0) es normal y se debe a la correlación TX/TX
        de la señal piloto OFDM consigo misma.
  \item Si aparecen picos en bins distintos de (0,0), puede indicar multipath
        en la interfaz loopback o un problema de configuración.
\end{itemize}
\end{infobox}

% ─────────────────────────────────────────────────────────────────────────────
\section{Cambios al código original de OpenISAC}
\label{sec:cambios}

Esta sección describe con detalle todos los cambios introducidos respecto al
repositorio original de OpenISAC en GitHub. Los cambios se dividen en:
(1) archivos nuevos, (2) archivos modificados, y (3) correcciones de bugs críticos
detectados durante el desarrollo de la extensión ZMQ SIL.

\subsection{Archivos nuevos}

{\small
\begin{tabularx}{\linewidth}{p{5.4cm}X}
\toprule
\textbf{Archivo} & \textbf{Descripción} \\
\midrule
\texttt{channel\_sim/src/}\newline\texttt{ChannelSimulator.cpp} &
  Canal simulado completo. Implementa el modelo físico descrito en la Sección~\ref{sec:canal},
  pipeline ZMQ, TUI ncurses interactiva y control UDP en tiempo real. \\
\addlinespace
\texttt{channel\_sim/src/}\newline\texttt{ChannelSimulator\_void.cpp} &
  Canal transparente. Pasa la señal sin modificar (referencia nula). \\
\addlinespace
\texttt{channel\_sim/include/}\newline\texttt{ChannelParams.hpp} &
  Parámetros del canal. Todos los campos son \texttt{std::atomic<>},
  permitiendo acceso concurrente seguro desde el hilo RT y la TUI/UDP. \\
\addlinespace
\texttt{channel\_sim/src/}\newline\texttt{ZmqInterface.cpp} &
  Interfaz ZeroMQ PUSH/PULL para el simulador de canal.
  Hace BIND en los puertos 5558 (recepción) y 5559 (transmisión). \\
\addlinespace
\texttt{channel\_sim/include/}\newline\texttt{ZmqInterface.hpp} &
  Cabecera de \texttt{ZmqInterface} y estructura \texttt{ZmqChannelSimConfig}. \\
\addlinespace
\texttt{scripts/run\_modulator.bash} &
  Watchdog: reinicia \texttt{OFDMModulator --zmq} automáticamente tras cualquier
  terminación no deseada. Parada limpia con Ctrl+C en el terminal
  del watchdog. \\
\bottomrule
\end{tabularx}
}

\subsection{Archivos modificados}

\subsubsection{src/OFDMModulator.cpp}

\paragraph{(a) Flag \texttt{--zmq} en la CLI}

El parser de argumentos de línea de comandos se extiende para reconocer el flag \texttt{--zmq}
que activa el modo SIL:

\begin{lstlisting}[style=cpp,
  caption={Detección del flag --zmq — OFDMModulator.cpp}]
bool zmq_mode = false;
// ...
for (int i = 1; i < argc; ++i) {
    if (std::string(argv[i]) == "--zmq") {
        zmq_mode = true;
    }
}
cfg.zmq_mode = zmq_mode;
if (zmq_mode) {
    // Usar YAML paired_frame_queue_size, con minimo 256
    // (la latencia ZMQ ~10-30 frames requiere cola ampliada)
    if (cfg.paired_frame_queue_size < 256) {
        cfg.paired_frame_queue_size = 256;
    }
}
\end{lstlisting}

\begin{infobox}[Nota sobre paired\_frame\_queue\_size]
En el código original, la cola de pareado TX/RX (\texttt{PairedFrameQueue}) tiene un
tamaño de 8 frames. El modo ZMQ introduce una latencia variable de 10–30 frames
entre TX y RX (latencia de la red loopback + procesado del ChannelSimulator). Con una
cola de 8 frames, el modulador descartaría frames TX constantemente. El mínimo de 256
frames proporciona $\approx 590$ ms de margen (a 434 Hz de frame rate), suficiente para
absorber ráfagas de latencia.
\end{infobox}

\paragraph{(b) Inicialización de sockets ZMQ}

En modo \texttt{--zmq}, la función \texttt{\_init\_usrp()} omite toda la inicialización
de hardware (USRP, streams) y en su lugar crea los sockets ZeroMQ:

\begin{lstlisting}[style=cpp,
  caption={Creación de sockets ZMQ en lugar de USRP — OFDMModulator.cpp}]
#ifdef OPENISAC_ZMQ_SUPPORT
if (_cfg.zmq_mode) {
    // Contexto ZMQ con 1 hilo I/O
    _zmq_ctx  = std::make_unique<zmq::context_t>(1);

    // Socket TX (PUSH): envia frames al ChannelSimulator (puerto 5558)
    // HWM=4: si el buffer del canal esta lleno, los sends bloquean 500ms y descartan
    _zmq_push = std::make_unique<zmq::socket_t>(*_zmq_ctx, zmq::socket_type::push);
    _zmq_push->set(zmq::sockopt::sndhwm, 4);
    _zmq_push->set(zmq::sockopt::sndtimeo, 500);  // 500ms timeout
    _zmq_push->connect("tcp://127.0.0.1:5558");

    // Socket RX sensing (PULL): recibe frames procesados (puerto 5559)
    // rcvtimeo=100ms: timeout corto para no bloquear el hilo indefinidamente
    _zmq_sensing_pull = std::make_unique<zmq::socket_t>(*_zmq_ctx, zmq::socket_type::pull);
    _zmq_sensing_pull->set(zmq::sockopt::rcvtimeo, 100);
    _zmq_sensing_pull->connect("tcp://127.0.0.1:5559");
    return;  // sin inicializacion de USRP
}
#endif
// ... inicializacion normal de USRP ...
\end{lstlisting}

\paragraph{(c) Hilo \_zmq\_sensing\_rx\_proc}

Nuevo hilo dedicado a recibir los frames del canal (puerto 5559) e inyectarlos
en el pipeline de sensing:

\begin{lstlisting}[style=cpp,
  caption={Hilo de recepción ZMQ sensing — OFDMModulator.cpp, lineas 1432--1482}]
void _zmq_sensing_rx_proc() {
    const size_t frame_size = _cfg.samples_per_frame();
    AlignedVector frame_buf(frame_size);
    size_t frame_received = 0;

    while (_running.load(std::memory_order_relaxed)) {
        zmq::message_t msg;
        try {
            auto res = _zmq_sensing_pull->recv(msg, zmq::recv_flags::none);
            if (!res || msg.size() == 0) continue;
        } catch (const zmq::error_t& e) {
            if (e.num() == ETERM) break;   // contexto destruido — salida limpia
            continue;
        }

        // Ensamblar frames OFDM completos a partir de mensajes ZMQ de tamano variable
        const size_t msg_samps = msg.size() / sizeof(std::complex<float>);
        const auto* src = static_cast<const std::complex<float>*>(msg.data());
        size_t src_pos = 0;

        while (src_pos < msg_samps) {
            const size_t to_copy = std::min(frame_size - frame_received,
                                            msg_samps - src_pos);
            std::memcpy(frame_buf.data() + frame_received, src + src_pos,
                        to_copy * sizeof(std::complex<float>));
            frame_received += to_copy;
            src_pos        += to_copy;

            if (frame_received >= frame_size) {
                const uint64_t seq = _zmq_rx_frame_seq.fetch_add(1);
                for (auto& ch : _sensing_channels) {
                    // Adquirir buffer del pool (NO asignar heap nuevo por frame):
                    // evita fuga de memoria de 14 GB en 37 s (ver Seccion 6.3)
                    AlignedVector ch_buf = ch->acquire_rx_frame();
                    if (ch_buf.size() != frame_size) ch_buf.resize(frame_size);
                    std::copy(frame_buf.begin(), frame_buf.end(), ch_buf.begin());
                    ch->push_zmq_rx_frame(std::move(ch_buf), seq);
                    // Si push falla (cola llena), push_zmq_rx_frame
                    // libera el buffer al pool automaticamente.
                }
                frame_received = 0;
            }
        }
    }
}
\end{lstlisting}

\paragraph{(d) Hilo \_tx\_proc — modo ZMQ}

El hilo transmisor se adapta para enviar los frames via ZMQ en lugar de al stream USRP.
Se añade una pausa temporal artificial para mantener la cadencia de frames:

\begin{lstlisting}[style=cpp,
  caption={Envío ZMQ y pacing temporal — OFDMModulator.cpp}]
// Envio del frame via ZMQ PUSH (en lugar de uhd::tx_stream::send())
const bool ok = _zmq_push->send(msg, zmq::send_flags::none).has_value();
// ...

// Guardia de underflow: solo activa el reinicio si NO estamos en modo ZMQ.
// En ZMQ, un fallo de send es un timeout de red (HWM), no un underflow HW.
if (!_cfg.zmq_mode && sent < frame_to_send.samples.size()) {
    _tx_underflow_restart_requested.store(true);
}

// Pacing temporal: dormir el tiempo equivalente a un frame
// para no saturar el canal ZMQ ni generar latencia acumulada.
// T_frame = samples_per_frame / sample_rate [segundos]
if (_cfg.zmq_mode) {
    std::this_thread::sleep_for(std::chrono::duration<double>(
        static_cast<double>(_cfg.samples_per_frame()) / _cfg.sample_rate));
}
\end{lstlisting}

\paragraph{(e) Registro de SIGTERM}

El código original solo registraba el manejador de señal para SIGINT (Ctrl+C).
SIGTERM (señal de terminación estándar de Linux/systemd) usaba el handler por
defecto, que termina el proceso \textbf{sin ejecutar ningún código de limpieza} y
sin imprimir mensajes de error:

\begin{lstlisting}[style=cpp,
  caption={Registro de SIGTERM — OFDMModulator.cpp}]
// Codigo original (solo SIGINT):
std::signal(SIGINT, &signal_handler);

// Codigo corregido (SIGINT + SIGTERM):
std::signal(SIGINT,  &signal_handler);
std::signal(SIGTERM, &signal_handler);   // terminar limpiamente con kill/systemd
\end{lstlisting}

\subsubsection{src/SensingChannel.cpp / include/SensingChannel.hpp}

\paragraph{(a) Método start\_zmq\_rx()}

Nuevo método público que arranca el hilo de procesado de sensing en modo ZMQ,
sin inicializar el hardware RX de la USRP:

\begin{lstlisting}[style=cpp,
  caption={start\_zmq\_rx() — SensingChannel.cpp, líneas 891--903}]
void SensingChannel::start_zmq_rx() {
    _stop_requested.store(false);
    _compute.bistatic_active = false;
    _compute.next_delay_estimation_frame_seq = 0;
    _rx_io.stream_start_time = uhd::time_spec_t(0.0);
    _rx_io.next_rx_frame_seq = 0;
    // Activar modo ZMQ loopback en la cola de pareado:
    // pop_pair() usara timeout en lugar de bloqueo infinito
    // para que el hilo de sensing no quede bloqueado si se descartan frames TX
    _rx_io.paired_queue->set_zmq_loopback(true);
    _compute.sensing_sender.start();
    // Sin hilo _rx_io.rx_thread: los frames los provee _zmq_sensing_rx_proc
    _compute.sensing_thread = std::thread(&SensingChannel::_sensing_loop, this);
}
\end{lstlisting}

\paragraph{(b) Métodos acquire\_rx\_frame() / release\_rx\_frame()}

Nuevos métodos para exponer el pool interno de buffers RX al hilo ZMQ externo,
necesarios para resolver la fuga de memoria descrita en la Sección~\ref{sec:bugs}:

\begin{lstlisting}[style=cpp,
  caption={Nuevos métodos de pool — SensingChannel.cpp / SensingChannel.hpp}]
// Adquirir buffer preasignado del pool (sin asignacion de heap)
AlignedVector SensingChannel::acquire_rx_frame() {
    return _rx_io.rx_frame_pool.acquire();
}

// Devolver buffer al pool (sin liberar memoria)
void SensingChannel::release_rx_frame(AlignedVector&& buf) {
    _rx_io.rx_frame_pool.release(std::move(buf));
}
\end{lstlisting}

\paragraph{(c) PairedFrameQueue: modo ZMQ loopback}

La clase interna \texttt{PairedFrameQueue} se extiende con un mecanismo de
\textbf{timeout en el lado TX} para el modo ZMQ. Sin este mecanismo, si la cola TX
se llena y se descartan frames TX, el hilo de sensing se bloquea indefinidamente
esperando el frame TX descartado (que nunca llegará):

\begin{lstlisting}[style=cpp,
  caption={Timeout ZMQ en PairedFrameQueue — SensingChannel.cpp, líneas 370--412}]
// Nueva bandera atomica: activada por set_zmq_loopback(true)
std::atomic<bool> _zmq_loopback{false};
void set_zmq_loopback(bool v) { _zmq_loopback.store(v); }

// Variante con timeout de 200 ms para modo ZMQ
bool _pop_next_tx_zmq(TxSymbolsFrame& out, bool& timed_out) {
    timed_out = false;
    using Clock = std::chrono::steady_clock;
    const auto deadline = Clock::now() + std::chrono::milliseconds(200);
    SPSCBackoff backoff;
    while (Clock::now() < deadline) {
        if (_tx_q.try_pop(out)) return true;
        if (_closed.load()) return false;  // cola cerrada
        backoff.pause();
    }
    timed_out = true;
    return false;   // timeout: frame TX descartado, resincronizan
}

// En pop_pair(): si hay timeout, descartar el frame RX y reintentar
if (_zmq_loopback.load()) {
    bool timed_out = false;
    if (!_pop_next_tx_zmq(tx_item, timed_out)) {
        _recycle_one(std::move(rx_item.samples));
        if (timed_out) { have_rx = false; continue; }  // descartar RX, reiniciar
        return false;  // cola cerrada
    }
}
\end{lstlisting}

\paragraph{(d) push\_zmq\_rx\_frame() — liberación al pool en caso de fallo}

Cuando la cola RX está llena y el push falla, el buffer se devuelve al pool
(en lugar de descartarse):

\begin{lstlisting}[style=cpp,
  caption={push\_zmq\_rx\_frame con liberación al pool — SensingChannel.cpp, líneas 905--916}]
bool SensingChannel::push_zmq_rx_frame(AlignedVector&& samples, uint64_t frame_seq) {
    RxSymbolsFrame rx_item;
    rx_item.samples = std::move(samples);
    rx_item.frame_seq = frame_seq;
    if (!_rx_io.paired_queue->push_rx(std::move(rx_item))) {
        LOG_RT_WARN_HZ(5) << "[Sensing CH " << _rx_io.logical_id
                          << "] ZMQ RX queue full, dropping frame " << frame_seq;
        // SPSCRingBuffer::try_push NO mueve el valor si la cola esta llena,
        // por lo que rx_item.samples es aun valido aqui.
        // Devolver el buffer al pool para mantener el ciclo acquire/release balanceado.
        _rx_io.rx_frame_pool.release(std::move(rx_item.samples));
        return false;
    }
    _rx_io.next_rx_frame_seq = frame_seq + 1;
    return true;
}
\end{lstlisting}

\subsubsection{CMakeLists.txt}

Se añaden dos nuevos targets de compilación condicionales al bloque
\texttt{if(CHANNEL\_SIM\_}\allowbreak\texttt{ZMQ\_FOUND)}:

\begin{lstlisting}[style=bash,
  caption={Nuevos targets en CMakeLists.txt}]
# Target ChannelSimulator (canal fisico con TUI)
add_executable(ChannelSimulator
    channel_sim/src/ChannelSimulator.cpp
    channel_sim/src/ZmqInterface.cpp
    src/AsyncLogger.cpp)
target_link_libraries(ChannelSimulator PRIVATE
    ${ZMQ_LIBRARIES} ncurses yaml-cpp uhd fftw3f ...)

# Target ChannelSimulator_void (canal transparente, sin ncurses)
add_executable(ChannelSimulator_void
    channel_sim/src/ChannelSimulator_void.cpp
    channel_sim/src/ZmqInterface.cpp
    src/AsyncLogger.cpp)
target_link_libraries(ChannelSimulator_void PRIVATE
    ${ZMQ_LIBRARIES} yaml-cpp uhd ...)
\end{lstlisting}

La variable de compilación \texttt{OPENISAC\_ZMQ\_SUPPORT} se define cuando ZMQ está
disponible y se propaga a \texttt{OFDMModulator} y \texttt{OFDMDemodulator}, activando
todas las secciones \texttt{\#ifdef OPENISAC\_ZMQ\_SUPPORT}.

\subsection{Correcciones de bugs críticos detectados}
\label{sec:bugs}

Durante el desarrollo de la extensión ZMQ SIL se identificaron y corrigieron tres bugs
críticos presentes en el código base original:

\subsubsection{Bug 1: Bloqueo infinito en PairedFrameQueue (modo ZMQ)}

\textbf{Descripción}: En modo ZMQ, cuando la cola TX (\texttt{PairedFrameQueue})
se llenaba y los frames TX se descartaban, el hilo de sensing llamaba a
\texttt{\_pop\_next\_tx()}, que gira indefinidamente esperando el siguiente frame TX.
El frame descartado nunca llegaría, bloqueando el hilo para siempre.

\textbf{Causa raíz}: \texttt{\_pop\_next\_tx()} no tiene mecanismo de timeout; fue
diseñado para uso con hardware donde los frames TX nunca se descartan.

\textbf{Solución}: Añadir \texttt{\_pop\_next\_tx\_zmq()} con timeout de 200 ms y
activarlo con el flag \texttt{set\_zmq\_loopback(true)} en \texttt{start\_zmq\_rx()}.
Véase el código en la Sección anterior.

\subsubsection{Bug 2: Fuga de memoria de 14,7 GB en \_zmq\_sensing\_rx\_proc}

\textbf{Descripción}: El proceso \texttt{OFDMModulator --zmq} terminaba
involuntariamente $\approx 37$ segundos después del arranque, eliminado por el
OOM killer del sistema operativo.

\textbf{Causa raíz}: El hilo \texttt{\_zmq\_sensing\_rx\_proc} asignaba un nuevo
\texttt{AlignedVector} de $\approx 921$ KB en cada frame recibido. El hilo de
sensing lo procesaba y lo \textbf{liberaba al pool} mediante
\texttt{rx\_frame\_pool.release()}. Sin embargo, el hilo ZMQ nunca adquiría del pool:
solo liberaba. El pool crecía sin límite:

\begin{equation}
  16.000 \text{ frames} \times 921 \text{ KB/frame} \approx 14{,}7 \text{ GB}
  \label{eq:mem_leak}
\end{equation}

A $\approx 434$ Hz de frame rate, la acumulación alcanza el límite de RAM en
$\approx 37$ segundos (con 16 GB de RAM instalada).

\textbf{Solución}: El hilo ZMQ ahora llama a \texttt{ch->acquire\_rx\_frame()} para
obtener buffers preasignados del pool. El hilo de sensing los devuelve con
\texttt{rx\_frame\_pool.release()}. El pool permanece en estado de equilibrio estacionario.

\subsubsection{Bug 3: SIGTERM sin manejador registrado}

\textbf{Descripción}: El proceso \texttt{OFDMModulator} terminaba silenciosamente
(sin mensajes de error) cuando recibía una señal SIGTERM (p.ej. de \texttt{kill},
\texttt{systemd} o \texttt{taskset}).

\textbf{Causa raíz}: El código original solo registraba el handler de señal para
SIGINT (Ctrl+C). SIGTERM usaba la acción por defecto del kernel, que termina el
proceso inmediatamente sin ejecutar destructores ni código de limpieza.

\textbf{Solución}: Añadir \texttt{std::signal(SIGTERM, \&signal\_handler)} junto a
la ya existente línea de SIGINT.

% ─────────────────────────────────────────────────────────────────────────────
\section{Manual de instalación}

\subsection{Dependencias del sistema}

\begin{lstlisting}[style=bash, caption={Instalación de dependencias en Ubuntu 22.04+}]
# Dependencias base de OpenISAC (igual que el repositorio original)
sudo apt update
sudo apt install -y \
    cmake build-essential pkg-config \
    libuhd-dev uhd-host \
    libboost-all-dev \
    libfftw3-dev \
    libyaml-cpp-dev \
    libncurses-dev

# ZeroMQ (necesario para el modo ZMQ SIL)
sudo apt install -y libzmq3-dev libczmq-dev

# cppzmq (cabeceras C++ de ZMQ — si no está en el sistema)
sudo apt install -y cppzmq-dev
# Alternativa: descargar zmq.hpp de https://github.com/zeromq/cppzmq

# AFF3CT (codec LDPC) — compilar desde fuentes
git clone --recursive https://github.com/aff3ct/aff3ct.git
cd aff3ct && mkdir build && cd build
cmake .. -DCMAKE_BUILD_TYPE=Release \
         -DAFF3CT_COMPILE_SHARED_LIB=ON \
         -DAFF3CT_COMPILE_STATIC_LIB=OFF
make -j$(nproc)
sudo make install && sudo ldconfig
cd ../..
\end{lstlisting}

\subsection{Compilación}

\begin{lstlisting}[style=bash, caption={Compilación del proyecto}]
cd fork_OpenISAC

# Configurar (Release con optimizaciones O2)
cmake -B build -DCMAKE_BUILD_TYPE=Release

# Compilar todos los targets
cmake --build build -j$(nproc)

# Binarios generados en build/:
#   OFDMModulator          — transmisor/receptor ISAC
#   OFDMDemodulator        — demodulador (modo bistico / UE)
#   ChannelSimulator       — canal simulado con TUI (requiere ZMQ)
#   ChannelSimulator_void  — canal transparente   (requiere ZMQ)
\end{lstlisting}

Si ZMQ no está disponible en el sistema, \texttt{cmake} mostrará un aviso y
omitirá los targets de canal simulado:
\begin{lstlisting}[style=plain]
-- ZeroMQ (libzmq + zmq.hpp) not found — ChannelSimulator will NOT be built.
\end{lstlisting}

\subsection{Configuración YAML}

Copiar una de las plantillas de configuración al directorio \texttt{build/}:

\begin{lstlisting}[style=bash]
cp config/Modulator_B210.yaml build/Modulator.yaml
# Verificar los parametros criticos:
#   sample_rate: 50000000      (Hz) — frecuencia de muestreo
#   center_freq: 2400000000    (Hz) — frecuencia portadora
#   fft_size: 1024             — numero de subportadoras
#   num_symbols: 100           — simbolos OFDM por frame
#   cp_length: 128             — longitud del prefijo ciclico
\end{lstlisting}

Para el modo ZMQ SIL, el bloque \texttt{channel\_sim} del YAML debe estar
\textbf{descomentado} (controla los puertos ZMQ y el puerto de control UDP):

\begin{lstlisting}[style=plain, caption={Bloque channel\_sim en Modulator.yaml}]
channel_sim:
  control_port: 9998       # Puerto UDP para comandos de control remoto
  buffer_size: 4096        # Tamano de buffer ZMQ (muestras)
  zmq:
    rx_bind_addr: "tcp://*:5558"   # ChannelSim recibe TX en este puerto
    tx_bind_addr: "tcp://*:5559"   # ChannelSim envia RX en este puerto
\end{lstlisting}

% ─────────────────────────────────────────────────────────────────────────────
\section{Manual de uso paso a paso}

\subsection{Modo 1: Hardware real (USRP)}
\label{sec:modo_hardware}

Este es el modo de operación original de OpenISAC. Requiere hardware USRP.

\begin{warnbox}[Prerrequisito hardware]
Conectar el USRP al PC via USB3 (B210) o SFP/PCIe (X310). Verificar la conexión:
\texttt{uhd\_find\_devices}. El USRP debe aparecer en la salida.
\end{warnbox}

\subsubsection{Terminal 1 — Configuración del sistema (una vez por sesión)}

\begin{lstlisting}[style=bash]
# Governor de CPU a performance (reduce jitter de frecuencia)
sudo ./scripts/set_performance.bash

# Aislar 6 cores para la aplicacion (OS usa solo cores 6-7)
# Sustituir 0-5 por el rango adecuado para el procesador del sistema
sudo ./scripts/isolate_cpus.bash 0-5
\end{lstlisting}

\subsubsection{Terminal 2 — OFDMModulator (sin flag --zmq)}

\begin{lstlisting}[style=bash]
cd build
# Lanzar con afinidad de CPU y prioridad RT:
sudo ../scripts/isolate_cpus.bash run ./OFDMModulator Modulator.yaml
\end{lstlisting}

El modulador se conecta al USRP, sintoniza la frecuencia y comienza a transmitir
frames OFDM. Los datos de sensing se envían por UDP al puerto 8888.

\subsubsection{Terminal 3 — Visualización}

\begin{lstlisting}[style=bash]
cd build
python3 ../scripts/plot_sensing_fast.py
\end{lstlisting}

Se abre la ventana de visualización con el mapa Rango-Doppler en tiempo real.
En ausencia de objetos, solo debería aparecer el bin DC (señal directa).

\subsection{Modo 2: Canal transparente (ChannelSimulator\_void)}
\label{sec:modo_void}

Este modo ejecuta la cadena ISAC completa sin hardware USRP y sin efectos de canal.
Es útil para verificar el funcionamiento del pipeline y establecer una referencia nula.

\begin{warnbox}[Orden de arranque obligatorio]
El ChannelSimulator (o \_void) debe arrancar \textbf{antes} que el OFDMModulator.
El simulador hace \texttt{bind()} en los puertos ZMQ; el modulador hace \texttt{connect()}.
Si se invierten, el modulador no encontrará el servidor ZMQ.
\end{warnbox}

\subsubsection{Terminal 1 — ChannelSimulator\_void}

\begin{lstlisting}[style=bash]
cd build
./ChannelSimulator_void Modulator.yaml
\end{lstlisting}

Salida esperada:
\begin{lstlisting}[style=plain]
[VOID] Canal transparente iniciado
[VOID]   RX (PULL): tcp://*:5558
[VOID]   TX (PUSH): tcp://*:5559
[VOID]   Buffer:    115200 muestras
[VOID] Ctrl+C para salir.

[VOID] Tramas retransmitidas: 1000
[VOID] Tramas retransmitidas: 2000
...
\end{lstlisting}

\subsubsection{Terminal 2 — OFDMModulator en modo ZMQ}

\begin{lstlisting}[style=bash]
cd build
./OFDMModulator --zmq Modulator.yaml
\end{lstlisting}

\subsubsection{Terminal 3 — Visualización}

\begin{lstlisting}[style=bash]
cd build
python3 ../scripts/plot_sensing_fast.py
\end{lstlisting}

\begin{infobox}[Resultado esperado — canal void]
El mapa Rango-Doppler \textbf{no debe mostrar picos de alta amplitud}
fuera del bin (rango=0, Doppler=0). Este resultado confirma que:
\begin{enumerate}
  \item El pipeline ZMQ funciona correctamente de extremo a extremo.
  \item No existen artefactos introducidos por el software de simulación.
  \item Cualquier pico que aparezca en los experimentos con canal simulado
        es atribuible exclusivamente al modelo de canal.
\end{enumerate}
\end{infobox}

\subsubsection{Parada limpia}

\begin{enumerate}
  \item Ctrl+C en Terminal 2 (OFDMModulator).
  \item Ctrl+C en Terminal 1 (ChannelSimulator\_void).
\end{enumerate}

\subsection{Modo 3: Canal simulado con efectos físicos (ChannelSimulator)}
\label{sec:modo_channel}

Este modo simula un canal radar monostático completo con los efectos físicos
descritos en la Sección~\ref{sec:canal}. Los parámetros se controlan en tiempo
real mediante la TUI ncurses.

\subsubsection{Terminal 1 — ChannelSimulator (con TUI)}

\begin{lstlisting}[style=bash]
cd build
./ChannelSimulator Modulator.yaml
\end{lstlisting}

Se abre la interfaz TUI ncurses. Para prioridad RT (hilo de señal en SCHED\_FIFO/80):

\begin{lstlisting}[style=bash]
sudo ./ChannelSimulator Modulator.yaml
\end{lstlisting}

\subsubsection{Terminal 2 — OFDMModulator en modo ZMQ}

\begin{lstlisting}[style=bash]
cd build
./OFDMModulator --zmq Modulator.yaml
# Para watchdog automatico (recomendado para sesiones largas):
sudo ../scripts/run_modulator.bash
\end{lstlisting}

\subsubsection{Terminal 3 — Visualización}

\begin{lstlisting}[style=bash]
cd build
python3 ../scripts/plot_sensing_fast.py
\end{lstlisting}

\subsubsection{Uso de la TUI ncurses}

\begin{center}
\begin{lstlisting}[style=plain]
 CHANNEL SIMULATOR — PANEL DE CONTROL     SR:50MHz  Fc:2.400GHz
 Parametro               Valor              Paso
─────────────────────────────────────────────────────────────────
 SISTEMA
   Potencia TX          0.000 dB          1
   Ganancia RX          0.000 dB          1
   Ruido AWGN         -60.000 dB          5
   Tipo de antena      Omnidireccional
   Targets             1

 TARGET 1
 > Distancia           25.000 m           10
   Velocidad           10.000 m/s         5
   MicroDop. Amp        0.000 rad         0.01
   MicroDop. Freq      10.000 Hz          1
   RCS                  1.000 m2          0.5

─────────────────────────────────────────────────────────────────
 up/dn Navegar | +/- <- -> Valor | j/k Paso x0.5/x2 | Q Salir
\end{lstlisting}
\end{center}

\begin{center}
\begin{tabularx}{0.85\linewidth}{lX}
\toprule
\textbf{Tecla} & \textbf{Acción} \\
\midrule
$\uparrow$ / $\downarrow$ & Navegar entre parámetros \\
$\rightarrow$ / \texttt{+} / \texttt{=} & Incrementar valor en un paso \\
$\leftarrow$ / \texttt{-} & Decrementar valor en un paso \\
\texttt{j} & Paso $\div 2$ (ajuste más fino) \\
\texttt{k} & Paso $\times 2$ (ajuste más grueso) \\
\texttt{Q} & Salir limpiamente del ChannelSimulator \\
\bottomrule
\end{tabularx}
\end{center}

\subsubsection{Experimentos sugeridos}

\textbf{Experimento 1 — Target estático} (verificación de rango):
\begin{enumerate}
  \item Configurar: Targets=1, Distancia=100 m, Velocidad=0 m/s, RCS=1 m².
  \item En el mapa Rango-Doppler, debe aparecer un pico en el bin de rango
        correspondiente a $d = \lfloor 2 \times 100 / c \times f_s \rceil$
        y en el bin Doppler=0.
  \item Variar la distancia con las flechas: el pico debe desplazarse en el eje de rango.
\end{enumerate}

\textbf{Experimento 2 — Target con velocidad} (verificación Doppler):
\begin{enumerate}
  \item Configurar: Targets=1, Distancia=50 m, Velocidad=20 m/s, RCS=5 m².
  \item El pico debe aparecer desplazado en el eje Doppler:
        $f_d = 2 \times 20 \times f_c / c$ Hz.
  \item Cambiar el signo de la velocidad: el pico debe saltar al lado opuesto.
\end{enumerate}

\textbf{Experimento 3 — Micro-Doppler} (verificación de firma):
\begin{enumerate}
  \item Configurar: Targets=1, Dist=50 m, Vel=15 m/s, MicroDop.Amp=0.3 rad,
        MicroDop.Freq=25 Hz.
  \item En la dimensión Doppler, deben aparecer \textit{sidebands} simétricas
        a $\pm k \times 25$ Hz del Doppler principal.
  \item El número y amplitud de sidebands depende de la amplitud $A_m$
        (expansión en series de Bessel de la modulación PM).
\end{enumerate}

\textbf{Experimento 4 — Antena directiva} (verificación de beam pattern):
\begin{enumerate}
  \item Configurar: Tipo antena = Directiva, Ancho haz = 20°,
        Target1: Dist=50 m, Vel=20 m/s, Ángulo=0° (on-boresight).
  \item Registrar la amplitud del pico.
  \item Aumentar el ángulo a 10°: la amplitud debe caer a $\approx -3$ dB.
  \item Aumentar a 20°: la amplitud debe caer a $\approx -12$ dB.
\end{enumerate}

\textbf{Experimento 5 — Múltiples targets}:
\begin{enumerate}
  \item Configurar: Targets=3, Target1: 50 m / 20 m/s,
        Target2: 150 m / −10 m/s, Target3: 300 m / 5 m/s.
  \item Deben aparecer tres picos distintos en el mapa Rango-Doppler,
        cada uno en la celda (rango, Doppler) correspondiente.
\end{enumerate}

\textbf{Experimento 6 — Degradación por ruido}:
\begin{enumerate}
  \item Configurar un target visible. Subir \textit{Ruido AWGN}
        desde $-60$ dB hacia $0$ dB.
  \item El pico del target debe ir desapareciendo progresivamente en el suelo de ruido.
  \item La SNR en el mapa Rango-Doppler: $\text{SNR}_{RD} \approx
        A_\text{eco}^2 / (2\sigma_n^2) \times N_r \times N_d$
        donde $N_r$ y $N_d$ son los tamaños de las FFT de rango y Doppler.
\end{enumerate}

\subsection{Control remoto por UDP (opcional)}

El \texttt{ChannelSimulator} acepta comandos en el puerto \texttt{control\_port} (9998
por defecto) para automatizar cambios de parámetros desde scripts:

\begin{lstlisting}[style=bash, caption={Ejemplos de control UDP desde Python}]
python3 - <<'EOF'
import socket, struct

s = socket.socket(socket.AF_INET, socket.SOCK_DGRAM)

# Cambiar velocidad del Target 1 a 30.5 m/s  (valor = vel * 100 = 3050)
s.sendto(b'VEL ' + struct.pack('>i', 3050), ('127.0.0.1', 9998))

# Cambiar distancia a 200.0 m  (valor = dist * 10 = 2000)
s.sendto(b'DIST' + struct.pack('>i', 2000), ('127.0.0.1', 9998))

# Activar micro-Doppler: amp=0.3 rad (valor = amp * 10000 = 3000)
s.sendto(b'MDAF' + struct.pack('>i', 3000), ('127.0.0.1', 9998))

# Frecuencia micro-Doppler: 25 Hz (valor = freq * 100 = 2500)
s.sendto(b'MDFF' + struct.pack('>i', 2500), ('127.0.0.1', 9998))
EOF
\end{lstlisting}

% ─────────────────────────────────────────────────────────────────────────────
\section{Rendimiento y optimizaciones en tiempo real}

\subsection{Tabla de ruido precalculada}

El generador estándar \texttt{std::normal\_distribution} con \texttt{mt19937} tiene
un coste de $\approx 15$ ns/muestra. Para un frame de $N_\text{sym} = 100$ símbolos
OFDM con $N_\text{FFT} + N_\text{CP} = 1024 + 128 = 1152$ muestras/símbolo, el frame
tiene $115\,200$ muestras × 2 (I/Q) = $230\,400$ valores Gaussianos:

\[
  T_\text{RNG} \approx 230\,400 \times 15 \text{ ns} = 3{,}5 \text{ ms}
\]

El periodo de frame a $f_s = 50$ MHz es:
\[
  T_\text{frame} = 115\,200 / 50 \times 10^6 = 2{,}3 \text{ ms}
\]

Por tanto, el generador estándar consumiría $>100\%$ del presupuesto de tiempo de
un frame, haciendo imposible la ejecución en tiempo real.

La tabla precalculada de $2^{17}$ muestras Gaussianas (512 KB) permite accesos a
$\approx 2$ ns/muestra (caché L3), reduciendo $T_\text{RNG} < 0{,}5$ ms, equivalente
al $\approx 20\%$ del presupuesto de frame.

\subsection{Fasor Doppler incremental}

La multiplicación compleja de rotación es $\approx 4$ operaciones FP (2 mul + 2 add),
con latencia $\approx 4$ ns en hardware x86 moderno. La alternativa (\texttt{cos}/\texttt{sin}
por muestra) tiene latencia $\approx 10-15$ ns. El ahorro total por frame (3 targets):
\[
  (15 - 4) \text{ ns} \times 115\,200 \text{ muestras} \times 3 \text{ targets}
  \approx 3{,}8 \text{ ms}
\]

\subsection{Hilo de señal con prioridad RT}

El hilo \texttt{signal\_loop} (canal de señal) se ejecuta con política
\texttt{SCHED\_FIFO} y prioridad 80 cuando se lanza con \texttt{sudo},
garantizando que no sea desalojado por procesos del sistema operativo
durante el procesado de un frame.

% ─────────────────────────────────────────────────────────────────────────────
\section{Resolución de problemas}

{\small
\begin{tabularx}{\linewidth}{p{4.8cm}X}
\toprule
\textbf{Síntoma} & \textbf{Causa / Solución} \\
\midrule
\texttt{Cannot load config: Modulator.yaml} &
  El binario busca el YAML en el directorio de trabajo.
  Ejecutar desde \texttt{build/} o pasar la ruta explícita:
  \texttt{./OFDMModulator --zmq ../build/Modulator.yaml} \\
\addlinespace
El modulador termina $\approx$ 37 s después de arrancar &
  Bug de fuga de memoria corregido en esta extensión.
  Recompilar desde el repositorio actualizado. \\
\addlinespace
Seq mismatch / drop de frames al inicio &
  El ChannelSimulator no ha arrancado antes que el modulador.
  Reiniciar en el orden correcto. \\
\addlinespace
Mapa Rango-Doppler sin pico de target &
  Verificar: \texttt{num\_targets $\geq$ 1} en la TUI,
  \texttt{tx\_power\_db} y \texttt{rcs} $> 0$,
  \texttt{noise\_power\_db} $\leq -40$ dB. \\
\addlinespace
TUI con caracteres extraños &
  Verificar locale UTF-8: \texttt{echo \$LANG}
  (debe mostrar \texttt{es\_ES.UTF-8} o \texttt{en\_US.UTF-8}). \\
\addlinespace
\texttt{[WARN] Sin prioridad RT} &
  Lanzar \texttt{ChannelSimulator} con \texttt{sudo} para
  activar \texttt{SCHED\_FIFO/80} en el hilo de señal. \\
\addlinespace
\texttt{ZMQ error: ETERM} &
  El contexto ZMQ se ha destruido (normal en apagado).
  No es un error: indica salida limpia del hilo de recepción. \\
\addlinespace
\texttt{OFDMModulator --zmq} termina nada más arrancar &
  El ChannelSimulator no ha arrancado aún. El socket TX
  (PUSH + HWM=4 + sndtimeo=500 ms) descarta frames si
  no hay receptor. Arrancar el canal primero. \\
\bottomrule
\end{tabularx}
}% end \small

% ─────────────────────────────────────────────────────────────────────────────
\section{Resumen de todos los cambios frente al repositorio original}

\begin{tabularx}{\linewidth}{p{5.5cm}lX}
\toprule
\textbf{Archivo} & \textbf{Tipo} & \textbf{Cambio} \\
\midrule
\nolinkurl{channel_sim/src/ChannelSimulator.cpp}      & Nuevo & Canal físico: retardo, Doppler, micro-Doppler, ec. radar, AWGN, TUI \\
\nolinkurl{channel_sim/src/ChannelSimulator_void.cpp} & Nuevo & Canal transparente (referencia nula) \\
\nolinkurl{channel_sim/include/ChannelParams.hpp}     & Nuevo & Parámetros del canal (atómicos, thread-safe) \\
\nolinkurl{channel_sim/src/ZmqInterface.cpp}          & Nuevo & Interfaz ZMQ PUSH/PULL (BIND en 5558/5559) \\
\nolinkurl{channel_sim/include/ZmqInterface.hpp}      & Nuevo & Cabecera de ZmqInterface y ZmqChannelSimConfig \\
\nolinkurl{scripts/run_modulator.bash}                & Nuevo & Watchdog: reinicio automático del modulador \\
\midrule
\nolinkurl{src/OFDMModulator.cpp} & Modificado &
  (1) Flag \texttt{--zmq} en CLI;
  (2) Inicialización ZMQ en lugar de USRP;
  (3) Hilo \texttt{\_zmq\_sensing\_rx\_proc};
  (4) Pacing temporal TX;
  (5) Guard de underflow modo ZMQ;
  (6) SIGTERM registrado;
  (7) Fix fuga de memoria: \texttt{acquire\_rx\_frame()} \\
\addlinespace
\nolinkurl{src/SensingChannel.cpp} & Modificado &
  (1) Método \texttt{start\_zmq\_rx()};
  (2) Método \texttt{push\_zmq\_rx\_frame()};
  (3) \texttt{acquire\_rx\_frame()} / \texttt{release\_rx\_frame()};
  (4) \texttt{PairedFrameQueue}: timeout ZMQ loopback;
  (5) Fix liberación al pool en push fallido \\
\addlinespace
\nolinkurl{include/SensingChannel.hpp} & Modificado &
  Declaraciones de \texttt{start\_zmq\_rx()}, \texttt{push\_zmq\_rx\_frame()},
  \texttt{acquire\_rx\_frame()}, \texttt{release\_rx\_frame()} \\
\addlinespace
\nolinkurl{CMakeLists.txt} & Modificado &
  Targets \texttt{ChannelSimulator} y \texttt{ChannelSimulator\_void};
  define \texttt{OPENISAC\_ZMQ\_SUPPORT} \\
\addlinespace
\nolinkurl{scripts/*.bash} & Modificado &
  Conversión CRLF $\rightarrow$ LF (scripts inejecutables en Linux original) \\
\bottomrule
\end{tabularx}

% ─────────────────────────────────────────────────────────────────────────────
\vfill
\begin{center}
\rule{0.6\linewidth}{0.4pt}\\[0.5em]
{\small\color{codegray}
  \textbf{OpenISAC ZMQ SIL Extension} —
  Basado en el proyecto OpenISAC original (GitHub).\\
  Documento de referencia para publicación en congreso ISAC.\\
  Generado con XeLaTeX · fork\_OpenISAC · Mayo 2026.
}
\end{center}

\end{document}
